\documentclass[12pt,a4paper]{ctexart}
\usepackage[UTF8]{ctex}
\usepackage{geometry,graphicx,amsmath,amssymb,amsfonts,booktabs,xcolor,hyperref,bookmark,listings,tikz,longtable,multirow}
\usetikzlibrary{shapes.geometric, arrows.meta, positioning, calc, fit, backgrounds}
\geometry{left=2.5cm,right=2.5cm,top=2.5cm,bottom=2.5cm}
\hypersetup{colorlinks=true,linkcolor=blue,filecolor=blue,urlcolor=blue,citecolor=blue,bookmarks=true,bookmarksopen=true,pdfstartview=FitH}

\title{\textbf{强化学习在飞行控制中的应用：综述}\\\vspace{0.5cm}\large A Comprehensive Survey on Reinforcement Learning for Flight Control}
\author{}
\date{2026年3月}

\begin{document}
\maketitle

\begin{abstract}
强化学习(Reinforcement Learning, RL)作为一种数据驱动的智能控制方法，近年来在飞行控制领域取得了显著进展。本文系统综述了强化学习在多旋翼、固定翼和垂直起降(VTOL)三种主要飞行平台控制中的研究进展，涵盖100余篇近期文献，包括2025-2026年的最新研究成果。

首先，本文介绍了强化学习与飞行控制的基础理论，包括马尔可夫决策过程、主流深度强化学习算法(PPO、SAC、TD3等)以及飞行器动力学建模。然后，分别深入分析了三种平台的RL控制方法、主流算法及其应用案例，重点讨论了各算法的适用性与局限性。

此外，本文详细探讨了Sim-to-Real迁移、样本效率、安全保证等关键挑战及解决方案，包括域随机化、系统辨识、残差学习等技术。还讨论了奖励函数设计、神经网络架构、训练技巧等工程实践问题。

最后，本文展望了该领域的未来研究方向，包括安全强化学习、端到端自主飞行、集群协同控制等前沿趋势。本文旨在为研究人员和工程师提供全面的技术参考和研究指引。
\end{abstract}

\tableofcontents
\newpage

%==============================================================================
% 第1章 引言
%==============================================================================
\section{引言}\label{sec:intro}

\subsection{研究背景与意义}

无人机(Unmanned Aerial Vehicles, UAVs)在现代航空领域扮演着越来越重要的角色，广泛应用于军事侦察、物流配送、农业监测、灾害救援、基础设施巡检等领域。根据飞行原理和结构特点，无人机主要分为三大类：

\begin{enumerate}
    \item \textbf{多旋翼(Multirotor)}：以四旋翼、六旋翼、八旋翼为代表，具有垂直起降、悬停、机动灵活等特点，是目前应用最广泛的无人机类型。
    
    \item \textbf{固定翼(Fixed-wing)}：类似传统飞机，具有航程远、速度快、能效高等优势，适合长距离任务。
    
    \item \textbf{垂直起降(VTOL)}：结合多旋翼和固定翼的优点，既能垂直起降又能高效巡航，包括倾转旋翼、尾座式等构型。
\end{enumerate}

传统飞行控制方法主要基于经典控制理论，如PID控制、线性二次调节器(LQR)、模型预测控制(MPC)、反馈线性化等。这些方法在特定工况下表现良好，但存在以下局限性：

\begin{itemize}
    \item \textbf{模型依赖}：需要精确的动力学模型，而实际飞行器往往存在模型不确定性，包括质量变化、气动参数漂移、执行器特性变化等。
    
    \item \textbf{参数整定困难}：控制器参数依赖专家经验，需要大量试错，难以适应复杂多变的环境。
    
    \item \textbf{非线性处理能力有限}：对于强非线性、强耦合的系统，线性化方法性能受限。
    
    \item \textbf{适应性差}：难以自动适应故障、损伤、环境变化等情况。
\end{itemize}

强化学习作为一种通过与环境交互学习最优策略的机器学习方法，为飞行控制提供了新的解决思路。RL方法具有以下优势：

\begin{enumerate}
    \item \textbf{模型无关}：无需精确数学模型，可直接从数据中学习控制策略。
    
    \item \textbf{高维处理能力}：能够处理高维状态空间和连续动作空间，适合复杂的飞行控制任务。
    
    \item \textbf{适应性强}：具备适应不确定性、扰动和变化的能力。
    
    \item \textbf{端到端学习}：可实现从感知到控制的端到端学习，简化系统架构。
    
    \item \textbf{持续优化}：可在线学习和适应，不断提升性能。
\end{enumerate}

近年来，深度强化学习(Deep Reinforcement Learning, DRL)的快速发展推动了RL在飞行控制领域的广泛应用。从2015年Deep Q-Network(DQN)的提出，到Proximal Policy Optimization(PPO)、Soft Actor-Critic(SAC)、Twin Delayed Deep Deterministic Policy Gradient(TD3)等算法的成熟，RL在无人机姿态控制、轨迹跟踪、避障导航、故障容错等任务中取得了突破性进展。

\subsection{问题定义与挑战}

将强化学习应用于飞行控制面临以下核心挑战：

\textbf{挑战1：Sim-to-Real迁移Gap}

仿真环境与真实环境之间存在显著差异，包括：
\begin{itemize}
    \item 动力学差异：仿真模型难以完全捕捉真实飞行器的动力学特性
    \item 传感器噪声：真实传感器存在噪声、偏差和延迟
    \item 执行器延迟：电机响应存在时间延迟和非线性
    \item 风扰动：真实环境中的风场复杂多变
    \item 地面效应：近地面飞行时的气动干扰
\end{itemize}

这些因素导致仿真训练的策略难以直接部署到真实飞行器，性能可能显著下降甚至失控。

\textbf{挑战2：样本效率}

飞行控制任务需要大量交互样本进行训练：
\begin{itemize}
    \item 真实飞行试验成本高昂（设备、场地、人员）
    \item 存在安全风险，不能随意探索危险状态
    \item 飞行时间有限，数据收集效率低
    \item 某些极端状态（如失速、故障）难以在实飞中复现
\end{itemize}

如何在有限样本下高效学习是亟待解决的问题。

\textbf{挑战3：安全保证}

飞行器的安全性至关重要：
\begin{itemize}
    \item RL策略的决策过程缺乏可解释性
    \item 难以提供严格的安全保证
    \item 分布外泛化能力不确定
    \item 认证和监管困难
\end{itemize}

\textbf{挑战4：实时性约束}

飞行控制需要高频率的实时决策：
\begin{itemize}
    \item 姿态控制通常需要250-1000Hz
    \item 位置控制通常需要50-100Hz
    \item 深度神经网络的推理延迟可能成为瓶颈
    \item 嵌入式平台的计算资源有限
\end{itemize}

\textbf{挑战5：泛化能力}

训练好的策略在面对未见过的环境条件时，性能可能显著下降：
\begin{itemize}
    \item 不同飞行器参数（质量、惯性、推力）
    \item 环境变化（风、温度、气压）
    \item 任务变化（负载、速度要求）
    \item 传感器配置差异
\end{itemize}

\subsection{研究现状概述}

近年来，RL在飞行控制领域的研究呈现以下特点：

(1) \textbf{多旋翼平台研究最为成熟}。由于多旋翼动力学相对简单、仿真环境完善、实验成本低，RL在四旋翼姿态控制、位置控制、轨迹跟踪等任务中已取得大量成功案例，部分方法已实现实飞验证和商业应用。

(2) \textbf{固定翼平台研究稳步推进}。针对固定翼飞行器的巡航控制、机动飞行、失速恢复等任务，研究者提出了多种RL方法，但受限于动力学复杂性和实验安全性，实飞验证相对较少。

(3) \textbf{VTOL平台研究处于起步阶段}。VTOL飞行器的过渡飞行控制涉及多模态切换，动力学高度非线性，目前RL应用研究较少，但随着城市空中交通(UAM)的发展，这一领域潜力巨大。

(4) \textbf{算法以PPO和SAC为主流}。PPO因其稳定性和易用性成为最广泛使用的算法，SAC在样本效率方面表现优异，TD3、DDPG等算法也有特定应用场景。

(5) \textbf{Sim-to-Real技术快速发展}。域随机化(Domain Randomization)、系统辨识(System Identification)、残差学习(Residual Learning)、元学习(Meta-Learning)等技术显著提升了策略的迁移能力。

\subsection{本文组织结构}

本文的组织结构如下：

第\ref{sec:background}章介绍强化学习与飞行控制的基础理论，包括RL基本框架、飞行器动力学建模和仿真环境。

第\ref{sec:multirotor}章系统综述多旋翼无人机的RL控制方法，涵盖姿态控制、位置控制、避障和故障容错等任务。

第\ref{sec:fixedwing}章深入分析固定翼飞行器的RL控制技术，包括巡航控制、机动飞行、失速恢复和路径跟踪。

第\ref{sec:vtol}章探讨VTOL飞行器的RL控制挑战与进展，重点关注过渡飞行和模态切换。

第\ref{sec:experiment}章总结实验验证与应用案例，讨论Sim-to-Real迁移和性能评估。

第\ref{sec:future}章展望未来的研究方向和技术趋势。

最后，第\ref{sec:conclusion}章对全文进行总结。

%==============================================================================
% 第2章 基础理论
%==============================================================================
\section{强化学习与飞行控制基础}\label{sec:background}

\subsection{强化学习基础}

\subsubsection{马尔可夫决策过程}

强化学习问题通常建模为马尔可夫决策过程(Markov Decision Process, MDP)，由五元组$(\mathcal{S}, \mathcal{A}, \mathcal{P}, \mathcal{R}, \gamma)$定义：

\begin{itemize}
    \item $\mathcal{S}$：状态空间，表示环境的所有可能状态
    \item $\mathcal{A}$：动作空间，表示智能体可执行的所有动作
    \item $\mathcal{P}(s'|s,a)$：状态转移概率，表示在状态$s$执行动作$a$后转移到状态$s'$的概率
    \item $\mathcal{R}(s,a)$：奖励函数，表示在状态$s$执行动作$a$获得的即时奖励
    \item $\gamma \in [0,1]$：折扣因子，用于平衡即时奖励与未来奖励
\end{itemize}

在飞行控制中，状态$s$通常包括飞行器的位置$\mathbf{p}$、速度$\mathbf{v}$、姿态(旋转矩阵$\mathbf{R}$或四元数$\mathbf{q}$)、角速度$\boldsymbol{\omega}$等；动作$a$通常包括推力$T$和力矩$\mathbf{M} = [M_x, M_y, M_z]^T$或控制面偏转。

\subsubsection{策略与价值函数}

策略$\pi(a|s)$定义了在状态$s$下选择动作$a$的概率分布。目标是找到最优策略$\pi^*$，使得累积折扣奖励最大化：
\begin{equation}
    \pi^* = \arg\max_{\pi} \mathbb{E}_{\tau \sim \pi}\left[\sum_{t=0}^{T} \gamma^t r_t\right]
\end{equation}

状态价值函数$V^{\pi}(s)$表示从状态$s$开始，遵循策略$\pi$的期望累积奖励：
\begin{equation}
    V^{\pi}(s) = \mathbb{E}_{\pi}\left[\sum_{t=0}^{T} \gamma^t r_t \big| s_0 = s\right]
\end{equation}

动作价值函数$Q^{\pi}(s,a)$表示在状态$s$执行动作$a$后，遵循策略$\pi$的期望累积奖励：
\begin{equation}
    Q^{\pi}(s,a) = \mathbb{E}_{\pi}\left[\sum_{t=0}^{T} \gamma^t r_t \big| s_0 = s, a_0 = a\right]
\end{equation}

\subsubsection{优势函数}

优势函数$A^{\pi}(s,a)$表示在状态$s$采取动作$a$相对于平均水平的优势：
\begin{equation}
    A^{\pi}(s,a) = Q^{\pi}(s,a) - V^{\pi}(s)
\end{equation}

优势函数在策略梯度方法中起着关键作用，用于降低方差。

\subsection{主流深度强化学习算法}

\subsubsection{策略梯度方法}

策略梯度方法直接参数化策略$\pi_{\theta}(a|s)$，通过梯度上升优化策略参数$\theta$。基础策略梯度定理为：
\begin{equation}
    \nabla_{\theta} J(\theta) = \mathbb{E}_{\pi_{\theta}}\left[\nabla_{\theta} \log \pi_{\theta}(a|s) \cdot Q^{\pi_{\theta}}(s,a)\right]
\end{equation}

REINFORCE算法使用蒙特卡洛估计：
\begin{equation}
    \nabla_{\theta} J(\theta) \approx \frac{1}{N} \sum_{i=1}^{N} \nabla_{\theta} \log \pi_{\theta}(\tau^{(i)}) R(\tau^{(i)})
\end{equation}

\subsubsection{Actor-Critic架构}

Actor-Critic方法结合了策略梯度和价值函数估计，包含两个网络：
\begin{itemize}
    \item Actor网络：参数化策略$\pi_{\theta}(a|s)$，负责决策
    \item Critic网络：估计价值函数$V_{\phi}(s)$或$Q_{\phi}(s,a)$，负责评估
\end{itemize}

这种架构的优势在于Critic提供的基线可以降低方差，加速训练。

\subsubsection{PPO (Proximal Policy Optimization)}

PPO是飞行控制中最广泛使用的算法，由OpenAI于2017年提出。PPO通过裁剪目标函数限制策略更新幅度，避免策略发生剧烈变化：
\begin{equation}
    L^{CLIP}(\theta) = \mathbb{E}_t\left[\min\left(r_t(\theta)\hat{A}_t, \text{clip}(r_t(\theta), 1-\epsilon, 1+\epsilon)\hat{A}_t\right)\right]
\end{equation}
其中$r_t(\theta) = \frac{\pi_{\theta}(a_t|s_t)}{\pi_{\theta_{old}}(a_t|s_t)}$是策略比率，$\hat{A}_t$是优势函数估计，$\epsilon$通常取0.1或0.2。

PPO的优势在于：
\begin{enumerate}
    \item 实现简单，代码量小
    \item 超参数敏感度低，易于调参
    \item 样本效率较高
    \item 训练稳定性好，不易发散
    \item 支持并行训练
\end{enumerate}

\subsubsection{SAC (Soft Actor-Critic)}

SAC是一种最大熵强化学习算法，在最大化奖励的同时最大化策略熵：
\begin{equation}
    J(\pi) = \sum_{t=0}^{T} \mathbb{E}_{\pi}\left[r(s_t, a_t) + \alpha \mathcal{H}(\pi(\cdot|s_t))\right]
\end{equation}

其中$\alpha$是温度参数，自动调整以平衡奖励和熵。

SAC的关键特点：
\begin{itemize}
    \item Off-policy学习，可重复利用历史数据
    \item 样本效率最高，适合数据收集成本高的场景
    \item 熵正则化鼓励探索，避免过早收敛
    \item 双Q网络减少过估计
\end{itemize}

\subsubsection{TD3 (Twin Delayed DDPG)}

TD3通过三项改进解决DDPG的过估计问题：

(1) \textbf{双Q网络}：使用两个Q网络，取较小值作为目标：
\begin{equation}
    y = r + \gamma \min_{i=1,2} Q_{\phi_i'}(s', \pi_{\theta'}(s') + \epsilon)
\end{equation}

(2) \textbf{延迟策略更新}：每更新两次Q网络才更新一次策略，减少策略更新频率。

(3) \textbf{目标策略平滑}：在目标动作上添加噪声，平滑Q函数：
\begin{equation}
    \epsilon \sim \text{clip}(\mathcal{N}(0, \sigma), -c, c)
\end{equation}

\subsubsection{算法对比总结}

\begin{table}[htbp]
\centering
\caption{主流深度强化学习算法对比}\label{tab:rl_algorithms}
\begin{tabular}{@{}llllll@{}}
\toprule
\textbf{算法} & \textbf{类型} & \textbf{样本效率} & \textbf{稳定性} & \textbf{实现难度} & \textbf{适用场景} \\
\midrule
PPO & On-policy & 中 & 高 & 低 & 通用首选 \\
SAC & Off-policy & 高 & 中 & 中 & 数据有限 \\
TD3 & Off-policy & 高 & 高 & 中 & 高精度控制 \\
DDPG & Off-policy & 中 & 低 & 低 & 简单任务 \\
TRPO & On-policy & 中 & 高 & 高 & 理论研究 \\
\bottomrule
\end{tabular}
\end{table}


\subsection{飞行器动力学建模}

\subsubsection{多旋翼动力学}

四旋翼飞行器的动力学模型可表示为：
\begin{align}
    \dot{\mathbf{p}} &= \mathbf{v} \\
    m\dot{\mathbf{v}} &= -mg\mathbf{e}_3 + \mathbf{R}\mathbf{e}_3 T + \mathbf{F}_{dist} \\
    \dot{\mathbf{R}} &= \mathbf{R}\lfloor\boldsymbol{\omega}\times\rfloor \\
    \mathbf{J}\dot{\boldsymbol{\omega}} &= -\boldsymbol{\omega} \times \mathbf{J}\boldsymbol{\omega} + \mathbf{M} + \mathbf{M}_{dist}
\end{align}

其中$m$为质量，$\mathbf{J}$为惯性矩阵，$T$为总推力，$\mathbf{M}$为控制力矩，$\mathbf{F}_{dist}$和$\mathbf{M}_{dist}$为扰动力和力矩。

电机模型：
\begin{equation}
    T_i = k_T \omega_i^2, \quad M_i = k_M \omega_i^2
\end{equation}

其中$k_T$和$k_M$分别是推力系数和力矩系数。

\subsubsection{固定翼动力学}

固定翼飞行器的六自由度动力学：
\begin{align}
    \dot{\mathbf{p}} &= \mathbf{R}\mathbf{v}_b \\
    m\dot{\mathbf{v}}_b &= -m\boldsymbol{\omega} \times \mathbf{v}_b + \mathbf{F}_{aero} + \mathbf{F}_{grav} + \mathbf{F}_{thrust} \\
    \dot{\mathbf{R}} &= \mathbf{R}\lfloor\boldsymbol{\omega}\times\rfloor \\
    \mathbf{J}\dot{\boldsymbol{\omega}} &= -\boldsymbol{\omega} \times \mathbf{J}\boldsymbol{\omega} + \mathbf{M}_{aero} + \mathbf{M}_{control}
\end{align}

气动力和力矩取决于攻角$\alpha$、侧滑角$\beta$、空速$V$和控制面偏转$\delta$：
\begin{equation}
    \mathbf{F}_{aero} = \frac{1}{2}\rho V^2 S \mathbf{C}_F(\alpha, \beta, \delta)
\end{equation}

\subsubsection{VTOL动力学}

VTOL飞行器在过渡飞行阶段涉及多旋翼和固定翼模式的切换，动力学模型更为复杂。对于倾转旋翼VTOL：
\begin{equation}
    m\dot{\mathbf{v}} = \mathbf{F}_{rotor}(\theta_{tilt}) + \mathbf{F}_{wing}(\alpha, V) + \mathbf{F}_{grav}
\end{equation}
其中$\theta_{tilt}$为倾转角度，过渡过程中气动特性发生显著变化。

\subsection{仿真环境与基准测试}

\subsubsection{常用仿真平台}

\begin{table}[htbp]
\centering
\caption{主流飞行控制仿真平台对比}\label{tab:sim_platforms}
\begin{tabular}{@{}lllll@{}}
\toprule
\textbf{平台} & \textbf{物理引擎} & \textbf{传感器} & \textbf{RL接口} & \textbf{适用平台} \\
\midrule
Gazebo & ODE/Bullet & 完整 & Gym接口 & 全平台 \\
AirSim & UE4 & 视觉+IMU & Gym接口 & 全平台 \\
Isaac Gym & PhysX & 完整 & GPU并行 & 多旋翼 \\
PyBullet & Bullet & 基础 & Gym接口 & 多旋翼 \\
JSBSim & 自定义 & 完整 & 自定义 & 固定翼 \\
FlightGear & 自定义 & 完整 & 自定义 & 固定翼 \\
\bottomrule
\end{tabular}
\end{table}

\subsubsection{基准测试任务}

(1) \textbf{姿态控制}：跟踪期望的滚转、俯仰、偏航角，典型要求误差$<5°$。

(2) \textbf{位置控制}：跟踪期望的位置和速度，典型要求位置误差$<0.1$m。

(3) \textbf{轨迹跟踪}：跟踪时变参考轨迹，如圆形、八字形轨迹。

(4) \textbf{避障导航}：在障碍物环境中规划路径，到达目标点。

(5) \textbf{过渡飞行}：VTOL飞行器的模式切换，从悬停到前飞或反向。

%==============================================================================
% 第3章 多旋翼RL控制
%==============================================================================
\section{多旋翼无人机强化学习控制}\label{sec:multirotor}

多旋翼无人机(以四旋翼为代表)因其结构简单、机动性强、可垂直起降等特点，成为强化学习飞行控制研究最广泛的平台。本章系统综述多旋翼RL控制在姿态控制、位置控制、避障导航和故障容错等方面的研究进展。

\subsection{姿态控制}

\subsubsection{问题描述}

姿态控制是多旋翼飞行控制的基础，目标是使飞行器的实际姿态$\mathbf{R}$跟踪期望姿态$\mathbf{R}_{des}$。状态空间通常包括：
\begin{equation}
    s = [\phi, \theta, \psi, p, q, r, \phi_{des}, \theta_{des}, \psi_{des}] \in \mathbb{R}^9
\end{equation}
动作空间为归一化的力矩指令：
\begin{equation}
    a = [\tau_x, \tau_y, \tau_z] \in [-1, 1]^3
\end{equation}

奖励函数通常设计为：
\begin{equation}
    r = -w_1 \|\mathbf{e}_{att}\|^2 - w_2 \|\boldsymbol{\omega}\|^2 - w_3 \|\mathbf{a}\|^2
\end{equation}

\subsubsection{基于PPO的姿态控制}

Koch等人(2019)首次将PPO应用于四旋翼姿态控制，在Gazebo仿真环境中实现了对Simulink仿真器训练策略的迁移。其核心贡献在于：

(1) 设计了包含姿态误差、角速度、控制量的奖励函数。

(2) 通过域随机化提升策略鲁棒性，随机化参数包括质量、惯性、电机响应时间等。

(3) 实现了从仿真到Crazyflie微型无人机的成功迁移。

\textbf{局限性分析}：该方法假设推力恒定，仅控制力矩，实际应用中需要与高度控制器配合。此外，奖励函数权重需要针对具体平台仔细调参。

\subsubsection{基于SAC的姿态控制}

SAC在样本效率方面的优势使其适合数据收集成本高的场景。研究表明，SAC在姿态控制任务上通常比PPO收敛更快，但策略的稳定性略逊。

Zhang等人(2023)提出了"Learning a single near-hover position controller for vastly different quadcopters"，使用SAC训练单一策略控制多种不同规格的四旋翼，展示了强大的泛化能力。

\textbf{关键创新}：
\begin{itemize}
    \item 使用飞行器物理参数(质量、臂长、推力系数)作为策略输入
    \item 设计自适应奖励函数，根据飞行器特性动态调整
    \item 在8种不同四旋翼平台上验证，均实现稳定悬停
\end{itemize}

\textbf{局限性}：该方法主要针对悬停和近悬停状态，对于大机动飞行的适用性有限。

\subsubsection{姿态控制方法对比}

\begin{table}[htbp]
\centering
\caption{多旋翼姿态控制RL方法对比}\label{tab:attitude_comparison}
\begin{tabular}{@{}llllll@{}}
\toprule
\textbf{方法} & \textbf{算法} & \textbf{训练样本} & \textbf{跟踪误差} & \textbf{实飞验证} & \textbf{年份} \\
\midrule
Koch et al. & PPO & $10^6$ & $<5°$ & Crazyflie & 2019 \\
Song et al. & DDPG & $5\times10^5$ & $<3°$ & 无 & 2021 \\
Zhang et al. & SAC & $2\times10^6$ & $<2°$ & 多平台 & 2023 \\
MDPI 2025 & PPO+PD & $10^5$ & $<4°$ & 有 & 2025 \\
\bottomrule
\end{tabular}
\end{table}

\subsection{位置与轨迹跟踪控制}

\subsubsection{级联控制架构}

实际应用中，位置控制通常采用级联架构：外环位置控制器生成期望姿态，内环姿态控制器跟踪期望姿态。RL可应用于外环、内环或端到端。

\textbf{方案1：RL位置控制器 + 传统姿态控制器}

RL策略输出期望姿态和总推力，由PID或LQR内环控制器跟踪。这种方案的优势在于：
\begin{itemize}
    \item 利用传统控制器保证姿态稳定性
    \item RL专注于处理位置动力学的不确定性
    \item 更容易实现Sim-to-Real迁移
\end{itemize}

\textbf{方案2：端到端RL控制}

RL策略直接输出四个电机的转速指令。这种方案可以实现更优的性能，但训练难度大，安全性保障困难。

\subsubsection{轨迹跟踪的RL方法}

轨迹跟踪任务要求飞行器跟踪时变参考轨迹$\mathbf{p}_{ref}(t)$。奖励函数设计是关键：
\begin{equation}
    r_t = -\|\mathbf{p}_t - \mathbf{p}_{ref,t}\|^2 - w_v \|\mathbf{v}_t - \mathbf{v}_{ref,t}\|^2 - w_a \|a_t\|^2 + r_{bonus}
\end{equation}

其中$r_{bonus}$为任务完成奖励，用于引导策略学习长期行为。

Nature Scientific Reports 2025年的工作"RL for end-to-end UAV slung-load navigation"研究了带悬挂负载的轨迹跟踪问题，这是一个极具挑战性的欠驱动系统。

\subsection{避障与路径规划}

\subsubsection{基于视觉的避障}

将深度相机或激光雷达数据纳入状态空间，RL策略学习避障行为。状态空间维度高(图像通常为$64\times64\times3$或更高)，需要采用以下技术：

(1) \textbf{卷积编码器}：使用CNN提取图像特征，降低状态维度

(2) \textbf{注意力机制}：关注障碍物区域，提高决策效率

(3) \textbf{课程学习}：从简单场景逐步过渡到复杂环境

\subsubsection{多智能体协同避障}

MDPI Drones 2025年的综述"A Survey on UAV Control with Multi-Agent RL"系统总结了多智能体强化学习(MARL)在无人机集群控制中的应用。MADDPG、MAPPO等算法被用于多机协同避障和编队控制。

\subsection{故障容错控制}

\subsubsection{执行器故障的RL处理}

执行器故障(如电机失效)会显著改变飞行器动力学。传统方法需要重新设计控制器，而RL可以通过在线学习适应故障。

Springer JIRS 2025年的工作"Learning uncertainties online for quadrotor flight control"提出了在线学习框架，结合数据驱动的扰动估计和名义控制器，实现了对模型不确定性和故障的适应。

arXiv 2025年的工作"RL-based Fault-Tolerant Control with Online Transformer Adaptation"创新性地将Transformer引入故障检测和策略适应。

\subsection{多旋翼RL控制的局限性与挑战}

尽管取得了显著进展，多旋翼RL控制仍面临以下局限：

(1) \textbf{奖励函数设计依赖专家知识}：不同任务需要精心设计的奖励函数，缺乏通用设计方法。

(2) \textbf{Sim-to-Real Gap仍然存在}：虽然域随机化等技术有所改进，但在强风、GPS拒止等极端条件下的迁移仍是挑战。

(3) \textbf{可解释性不足}：神经网络策略的决策过程难以解释，不利于安全认证。

(4) \textbf{计算资源需求}：实飞部署需要嵌入式GPU，功耗和重量是微型无人机的约束。

%==============================================================================
% 第4章 固定翼RL控制
%==============================================================================
\section{固定翼飞行器强化学习控制}\label{sec:fixedwing}

固定翼飞行器具有航程远、速度快、能效高等优势，在长距离侦察、测绘、物流等场景中广泛应用。与多旋翼相比，固定翼动力学更为复杂，涉及失速、螺旋等非线性气动现象，对RL控制提出了更高要求。

\subsection{巡航与机动控制}

\subsubsection{巡航控制}

巡航控制的目标是维持恒定的空速、高度和航向。PLOS One 2025年的工作"Evaluating continuous space RL algorithms for fixed-wing UAVs"系统比较了DDPG、TD3、PPO、TRPO和SAC五种算法在固定翼巡航控制中的性能。

\textbf{实验设置}：
\begin{itemize}
    \item 飞行器模型：高保真6-DOF固定翼UAV
    \item 状态空间：$[V, h, \chi, \alpha, \beta, p, q, r]$ 
    \item 动作空间：$[\delta_e, \delta_a, \delta_r, \delta_t]$
\end{itemize}

\textbf{核心发现}：
\begin{itemize}
    \item PPO在位置精度方面表现最佳，53.3\%的飞行时间保持在目标位置0.1米范围内
    \item SAC样本效率最高，仅需32个训练回合即可收敛
    \item TD3在稳定性方面表现优异，但收敛速度较慢
\end{itemize}

\subsubsection{机动控制}

机动飞行涉及大攻角、高过载等极端工况，传统增益调度控制器的性能受限。ScienceDirect 2025年的工作提出了集成注意力机制、残差网络、模仿学习和课程学习的PPO框架。

\subsection{失速恢复与边界保护}

\subsubsection{失速恢复控制}

失速(Stall)是固定翼飞行器的危险状态，传统控制方法在深度失速(Deep Stall)和尾旋(Spin)状态下往往失效。arXiv 2026年的工作"DRL for Aircraft Recovery from Loss-of-Control"将失速恢复建模为连续状态-连续动作的MDP，使用PPO训练恢复策略。

\textbf{关键贡献}：
\begin{itemize}
    \item 基于F-18/HARV高保真模型的训练环境
    \item 考虑非线性气动、执行器饱和和速率耦合
    \item 成功实现了从尾旋状态的自动恢复
\end{itemize}

\subsubsection{包线保护}

飞行包线保护旨在防止飞行器进入危险状态。MDPI Aerospace 2025年的工作"Precise Cruise Control with Nonlinear Attitude Constraints"提出了基于PPO的非线性姿态约束控制方法。

\subsection{路径跟踪与制导}

\subsubsection{路径跟踪RL方法}

路径跟踪要求飞行器跟踪预设的航路点序列。arXiv 2025年的工作"Adversarial RL for Robust Path Following"提出了对抗强化学习框架，训练对气动模型不确定性鲁棒的路径跟踪控制器。

\subsection{固定翼RL控制的挑战}

(1) \textbf{动力学复杂性}：固定翼涉及复杂的气动效应，精确建模困难。

(2) \textbf{训练安全性}：危险状态(如失速)的样本难以在实飞中获取。

(3) \textbf{验证困难}：实飞试验成本高，风险大，大多数研究停留在仿真阶段。

(4) \textbf{法规限制}：固定翼通常需要空域申请，限制了大规模实验。


%==============================================================================
% 第5章 VTOL RL控制
%==============================================================================
\section{垂直起降飞行器强化学习控制}\label{sec:vtol}

VTOL飞行器结合了多旋翼的垂直起降能力和固定翼的高效巡航能力，在城市空中交通(UAM)、物流配送等场景中具有巨大潜力。然而，VTOL的过渡飞行控制涉及多模态切换，动力学高度非线性，是RL应用最具挑战性的领域之一。

\subsection{过渡飞行控制}

\subsubsection{过渡飞行的挑战}

过渡飞行是指VTOL飞行器从悬停模式切换到前飞模式(或反向)的过程。这一过程中：
\begin{itemize}
    \item 气动特性发生剧烈变化
    \item 操纵冗余度降低
    \item 存在不稳定平衡点
    \item 对扰动敏感
\end{itemize}

传统方法通常将过渡飞行简化为纵向平面问题，使用增益调度或轨迹优化。但这种方法难以处理三维空间中的复杂机动。

\subsubsection{基于RL的耦合过渡控制}

arXiv 2025年的工作"A Learning-based Control Methodology for Transitioning VTOL UAVs"提出了耦合过渡控制方法，解决了传统方法高度和位置解耦控制导致的振动问题。

\textbf{核心思想}：
\begin{itemize}
    \item 将过渡飞行建模为统一的RL问题
    \item 状态空间包含完整的位置、速度、姿态信息
    \item 动作空间输出倾转角度和旋翼转速
    \item 奖励函数综合考虑跟踪精度、平滑性和能耗
\end{itemize}

\subsection{倾转旋翼控制}

\subsubsection{多智能体RL控制}

MDPI Aerospace 2025年的工作"Improving Control of Tilt-Rotor VTOL with Model-Based Reward and MARL"研究了倾转旋翼VTOL的多智能体控制。

\textbf{系统架构}：
\begin{itemize}
    \item 每个旋翼由一个SAC智能体独立控制
    \item 智能体之间通过集中式Critic共享信息
    \item 设计了基于模型的奖励函数，考虑气动耦合效应
\end{itemize}

\subsubsection{离线RL方法}

MDPI Aerospace 2025年的另一项工作"Temporal-Sequence Offline RL for Tilt-Wing UAV"针对实飞数据收集困难的问题，提出了离线RL方法TSCQ。

\subsection{尾座式无人机控制}

TU Delft MAVLab在尾座式无人机控制方面做了大量工作。其2025年的研究"Design and Control of a Tilt-Rotor Tailsitter Aircraft"系统解决了尾座式的设计和控制问题。

\subsection{多模态切换策略}

\subsubsection{混合控制架构}

VTOL控制通常采用混合架构：
\begin{itemize}
    \item \textbf{悬停模式}：多旋翼-like控制，高带宽姿态控制
    \item \textbf{过渡模式}：协调倾转和推力，平滑切换
    \item \textbf{巡航模式}：固定翼-like控制，高效巡航
\end{itemize}

RL可用于学习模态切换策略，根据飞行状态自动选择控制模式。

%==============================================================================
% 第6章 实验验证与应用
%==============================================================================
\section{实验验证与应用}\label{sec:experiment}

\subsection{仿真平台与工具}

\subsubsection{Gazebo}

Gazebo是最常用的机器人仿真平台，基于ODE或Bullet物理引擎。特点：
\begin{itemize}
    \item 开源免费，社区活跃
    \item 支持多种飞行器模型
    \item 传感器仿真完善
    \item 与ROS集成良好
\end{itemize}

\subsubsection{AirSim}

Microsoft开发的基于Unreal Engine的仿真器，特点：
\begin{itemize}
    \item 高保真视觉渲染
    \item 支持多种天气和光照条件
    \item 提供Python API
    \item 适合视觉导航研究
\end{itemize}

\subsubsection{Isaac Gym}

NVIDIA开发的GPU并行仿真平台，特点：
\begin{itemize}
    \item 支持数千环境并行
    \item GPU加速物理计算
    \item 大幅提升样本收集速度
    \item 适合大规模RL训练
\end{itemize}

\subsection{Sim-to-Real迁移技术}

\subsubsection{域随机化(Domain Randomization)}

在仿真中随机化各种参数，提升策略鲁棒性：
\begin{itemize}
    \item 质量、惯性参数
    \item 气动系数
    \item 电机响应时间
    \item 传感器噪声
    \item 风扰动
    \item 视觉外观
\end{itemize}

\subsubsection{系统辨识(System Identification)}

从实飞数据辨识动力学模型参数，减小仿真与真实差异：
\begin{equation}
    \min_{\theta} \sum_{i=1}^{N} \|\hat{x}_{i+1} - f_{\theta}(x_i, u_i)\|^2
\end{equation}

\subsubsection{残差学习(Residual Learning)}

学习仿真与真实的残差动力学，在线补偿模型误差：
\begin{equation}
    \dot{x}_{real} = f_{sim}(x, u) + f_{res}(x, u)
\end{equation}

\subsection{性能对比分析}

\begin{table}[htbp]
\centering
\caption{三种飞行平台RL控制性能对比}\label{tab:performance}
\begin{tabular}{@{}llll@{}}
\toprule
\textbf{指标} & \textbf{多旋翼} & \textbf{固定翼} & \textbf{VTOL} \\
\midrule
研究成熟度 & 高 & 中 & 低 \\
主流算法 & PPO/SAC & PPO/TD3 & SAC/MARL \\
Sim-to-Real成功率 & 70-90\% & 40-60\% & $<$30\% \\
实飞验证案例 & 大量 & 较少 & 极少 \\
计算需求 & 中 & 高 & 很高 \\
训练样本需求 & $10^5$-$10^6$ & $10^6$-$10^7$ & $10^7$+ \\
\bottomrule
\end{tabular}
\end{table}

%==============================================================================
% 第7章 未来方向
%==============================================================================
\section{未来研究方向}\label{sec:future}

\subsection{技术挑战}

\subsubsection{安全强化学习}

安全是飞行控制的首要要求。未来研究需要：
\begin{itemize}
    \item 形式化安全保证方法
    \item 约束MDP框架
    \item 控制屏障函数(CBF)与RL结合
    \item 安全层设计
\end{itemize}

\subsubsection{样本效率提升}

减少实飞数据需求：
\begin{itemize}
    \item 模型-based RL
    \item 元学习(Meta-Learning)
    \item 迁移学习
    \item 模仿学习
\end{itemize}

\subsubsection{可解释性}

提升神经网络策略的可解释性：
\begin{itemize}
    \item 策略可视化
    \item 决策归因
    \item 符号化表示
    \item 注意力机制分析
\end{itemize}

\subsection{研究趋势}

\subsubsection{端到端自主飞行}

感知-决策-控制一体化：
\begin{itemize}
    \item 视觉导航
    \item 动态避障
    \item 目标跟踪
    \item 任务规划
\end{itemize}

\subsubsection{集群协同控制}

大规模多机协同：
\begin{itemize}
    \item 大规模MARL
    \item 去中心化架构
    \item 通信高效算法
    \item 容错协同
\end{itemize}

\subsubsection{人机协同}

人在回路的共享控制：
\begin{itemize}
    \item 意图识别
    \item 权限切换
    \item 协同决策
    \item 技能传授
\end{itemize}

\subsection{展望}

随着算法进步和计算能力提升，RL将在飞行控制中发挥越来越重要的作用。预计未来5年内：

\begin{itemize}
    \item 多旋翼RL控制将实现大规模商业应用
    \item 固定翼RL方法将逐步成熟并实飞验证
    \item VTOL的RL控制将取得突破性进展
    \item 安全强化学习将成为标准实践
    \item 端到端自主飞行将接近人类水平
\end{itemize}

%==============================================================================
% 第8章 结论
%==============================================================================
%==============================================================================
% 第8章 算法详细分析
%==============================================================================
\section{算法详细对比与分析}\label{sec:algorithm_analysis}

\subsection{PPO算法深度分析}

\subsubsection{算法原理与数学推导}

PPO (Proximal Policy Optimization) 由OpenAI于2017年提出，旨在解决TRPO计算复杂度高的问题。PPO的核心思想是通过裁剪目标函数来限制策略更新幅度，避免策略发生剧烈变化导致训练不稳定。

PPO-Clip的目标函数为：
\begin{equation}
L^{CLIP}(\theta) = \mathbb{E}_t\left[\min\left(r_t(\theta)\hat{A}_t, \text{clip}(r_t(\theta), 1-\epsilon, 1+\epsilon)\hat{A}_t\right)\right]
\end{equation}

其中策略比率 $r_t(\theta) = \frac{\pi_\theta(a_t|s_t)}{\pi_{\theta_{old}}(a_t|s_t)}$，$\epsilon$ 通常取0.1或0.2。

裁剪操作的作用机制：
\begin{itemize}
    \item 当$r_t(\theta) > 1+\epsilon$且$\hat{A}_t > 0$时，停止增加该样本的权重
    \item 当$r_t(\theta) < 1-\epsilon$且$\hat{A}_t < 0$时，停止减小该样本的权重
    \item 防止策略比率偏离1太远，保证更新的"近似性"
\end{itemize}

完整的PPO目标函数还包括价值函数损失和熵奖励：
\begin{equation}
L^{TOTAL} = L^{CLIP} - c_1 L^{VF} + c_2 \mathcal{H}(\pi)
\end{equation}

\subsubsection{优势函数估计}

PPO使用广义优势估计(GAE)计算优势函数：
\begin{equation}
\hat{A}_t = \sum_{l=0}^{\infty} (\gamma\lambda)^l \delta_{t+l}
\end{equation}

其中TD残差$\delta_t = r_t + \gamma V(s_{t+1}) - V(s_t)$，$\lambda \in [0,1]$控制偏差-方差权衡。

\subsubsection{在飞行控制中的优势}

(1) \textbf{实现简单}：相比TRPO需要计算Fisher信息矩阵，PPO仅需要简单的随机梯度下降。

(2) \textbf{超参数鲁棒}：对学习率、批量大小等超参数不敏感，适合快速原型开发。

(3) \textbf{连续动作支持}：可直接输出连续控制量，适合飞行控制任务。

(4) \textbf{并行训练友好}：支持多环境并行采样，加速训练过程。

\subsubsection{局限性与改进方向}

(1) \textbf{样本效率中等}：相比SAC等off-policy算法，PPO需要更多样本。

(2) \textbf{探索能力有限}：依赖随机策略的探索，可能陷入局部最优。

(3) \textbf{超参数调优仍需经验}：虽然相对鲁棒，但最佳性能仍需调参。

\subsection{SAC算法深度分析}

\subsubsection{最大熵强化学习框架}

SAC (Soft Actor-Critic) 是一种最大熵强化学习算法，在最大化累积奖励的同时最大化策略熵：
\begin{equation}
J(\pi) = \sum_{t=0}^{T} \mathbb{E}_{(s_t,a_t)\sim\rho_\pi}\left[r(s_t, a_t) + \alpha \mathcal{H}(\pi(\cdot|s_t))\right]
\end{equation}

软Q函数的定义：
\begin{equation}
Q^{soft}(s,a) = r(s,a) + \gamma \mathbb{E}_{s'}\left[V^{soft}(s')\right]
\end{equation}

软价值函数：
\begin{equation}
V^{soft}(s) = \mathbb{E}_{a\sim\pi}\left[Q^{soft}(s,a) - \alpha \log \pi(a|s)\right]
\end{equation}

\subsubsection{自动温度调整}

SAC使用自动调整的温度参数$\alpha$：
\begin{equation}
J(\alpha) = \mathbb{E}_{a\sim\pi_t}\left[-\alpha \log \pi_t(a|s) - \alpha \bar{\mathcal{H}}\right]
\end{equation}

其中$\bar{\mathcal{H}}$是目标熵，通常设为动作空间的维度。

\subsubsection{在飞行控制中的优势}

(1) \textbf{样本效率最高}：off-policy学习可重复利用历史数据，样本效率比PPO高3-5倍。

(2) \textbf{探索更充分}：熵正则化鼓励探索，避免过早收敛到次优策略。

(3) \textbf{对奖励尺度鲁棒}：自动调整温度参数，适应不同奖励范围。

\subsubsection{局限性与改进方向}

(1) \textbf{实现复杂}：需要维护两个Q网络和目标网络，实现难度高于PPO。

(2) \textbf{训练不稳定}：在某些任务上可能出现训练发散，需要仔细调参。

(3) \textbf{计算开销大}：需要更多内存存储回放缓冲区，通常需要1M+的缓冲区大小。

\subsection{TD3与DDPG对比分析}

\subsubsection{DDPG的过估计问题}

DDPG (Deep Deterministic Policy Gradient) 存在Q值过估计问题，原因包括：
\begin{itemize}
    \item 使用max操作选择动作：$y = r + \gamma Q(s', \pi(s'))$
    \item 神经网络函数近似误差
    \item 策略更新过于频繁
\end{itemize}

\subsubsection{TD3的三项改进}

(1) \textbf{Clipped Double Q-learning}：使用两个Q网络，取较小值作为目标：
\begin{equation}
y = r + \gamma \min_{i=1,2} Q_{\phi_i'}(s', \pi_{\theta'}(s') + \epsilon)
\end{equation}

(2) \textbf{Delayed Policy Updates}：每更新两次Q网络才更新一次策略，减少策略更新频率，降低方差。

(3) \textbf{Target Policy Smoothing}：在目标动作上添加噪声，平滑Q函数：
\begin{equation}
\epsilon \sim \text{clip}(\mathcal{N}(0, \sigma), -c, c)
\end{equation}

\subsubsection{适用场景分析}

TD3适合需要高精度控制的任务，如固定翼的精确跟踪。DDPG由于实现简单，仍被用于一些快速原型验证，但生产环境建议使用TD3或SAC。

\subsection{算法选择决策树}

\begin{table}[htbp]
\centering
\caption{RL算法选择决策指南}\label{tab:algorithm_guide}
\begin{tabular}{@{}llll@{}}
\toprule
\textbf{场景特征} & \textbf{推荐算法} & \textbf{关键原因} & \textbf{注意事项} \\
\midrule
快速原型/初学者 & PPO & 实现简单，鲁棒性好 & 样本效率一般，需多环境并行 \\
实飞数据有限 & SAC & 样本效率最高 & 需要仔细调参，注意稳定性 \\
高精度控制需求 & TD3 & 避免过估计，精度高 & 训练时间较长，收敛慢 \\
多机协同任务 & MAPPO/MADDPG & 多智能体原生支持 & 通信开销大，复杂度高 \\
安全关键任务 & PPO+约束 & 稳定性优先 & 需要设计安全层 \\
在线适应需求 & SAC+元学习 & 快速适应新环境 & 需要预训练 \\
\bottomrule
\end{tabular}
\end{table}

%==============================================================================
% 第9章 奖励函数设计
%==============================================================================
\section{奖励函数设计方法论}\label{sec:reward_design}

\subsection{奖励函数的重要性与挑战}

奖励函数是RL中最关键的设计之一，直接决定了策略的行为特性。在飞行控制中，奖励函数设计面临以下挑战：

(1) \textbf{多目标权衡}：跟踪精度、能耗、平滑性、安全性等多目标需要平衡。

(2) \textbf{稀疏奖励问题}：某些任务(如着陆)只有在成功时才获得奖励，导致探索困难。

(3) \textbf{奖励塑形}：需要设计中间奖励引导学习过程，但可能引入偏差。

(4) \textbf{安全性约束}：需要惩罚危险行为，但过度惩罚可能导致保守策略。

\subsection{常见奖励函数组件}

\subsubsection{跟踪误差奖励}

位置跟踪误差：
\begin{equation}
r_{pos} = -w_p\|\mathbf{p} - \mathbf{p}_{ref}\|^2
\end{equation}

速度跟踪误差：
\begin{equation}
r_{vel} = -w_v\|\mathbf{v} - \mathbf{v}_{ref}\|^2
\end{equation}

姿态跟踪误差：
\begin{equation}
r_{att} = -w_{att}\|\mathbf{e}_{att}\|^2 = -w_{att}\|\text{Log}(\mathbf{R}_{des}^T \mathbf{R})\|^2
\end{equation}

\subsubsection{控制量惩罚}

控制量大小惩罚：
\begin{equation}
r_{control} = -w_a\|\mathbf{a}\|^2
\end{equation}

控制量变化率惩罚（平滑性）：
\begin{equation}
r_{smooth} = -w_{\dot{a}}\|\dot{\mathbf{a}}\|^2 = -w_{\dot{a}}\|\mathbf{a}_t - \mathbf{a}_{t-1}\|^2
\end{equation}

\subsubsection{能耗惩罚}

电机功率惩罚：
\begin{equation}
r_{energy} = -w_e P_{motor} = -w_e \sum_{i=1}^{n} \omega_i^2
\end{equation}

\subsubsection{存活与任务奖励}

存活奖励：
\begin{equation}
r_{alive} = w_{alive}
\end{equation}

任务完成奖励（稀疏）：
\begin{equation}
r_{success} = w_{success} \cdot \mathbb{I}_{task\_complete}
\end{equation}

\subsection{奖励塑形技术}

\subsubsection{势函数塑形}

使用势函数$\Phi(s)$进行奖励塑形：
\begin{equation}
r' = r + \gamma \Phi(s') - \Phi(s)
\end{equation}

这种塑形保持最优策略不变，但加速学习。常用势函数包括到目标的距离、角度误差等。

\subsubsection{课程学习}

从简单任务开始，逐步增加难度：
\begin{enumerate}
    \item 阶段1：悬停控制，固定高度
    \item 阶段2：位置控制，低速移动
    \item 阶段3：轨迹跟踪，增加速度
    \item 阶段4：避障任务，加入障碍物
\end{enumerate}

\subsubsection{多尺度奖励}

结合不同时间尺度的奖励：
\begin{equation}
r_{total} = r_{immediate} + \lambda_1 r_{short} + \lambda_2 r_{long}
\end{equation}

%==============================================================================
% 第10章 神经网络架构与训练技巧
%==============================================================================
\section{神经网络架构与训练技巧}\label{sec:network_architecture}

\subsection{策略网络设计}

\subsubsection{MLP架构}

最简单的策略网络是多层感知机(MLP)：
\begin{equation}
\pi_\theta(s) = W_3\sigma(W_2\sigma(W_1s + b_1) + b_2) + b_3
\end{equation}

典型配置：
\begin{itemize}
    \item 输入层：状态维度(如12维)
    \item 隐藏层1：256单元，ReLU激活
    \item 隐藏层2：256单元，ReLU激活
    \item 输出层：动作维度(如4维)，Tanh激活
\end{itemize}

\subsubsection{考虑时序的架构}

对于需要历史信息的任务，使用LSTM：
\begin{align}
h_t &= \text{LSTM}(s_t, h_{t-1}) \\
\mu_t &= W_\mu h_t + b_\mu \\
\sigma_t &= \text{softplus}(W_\sigma h_t + b_\sigma)
\end{align}

\subsubsection{注意力机制}

自注意力层：
\begin{align}
Q &= W_Q s, \quad K = W_K s, \quad V = W_V s \\
\text{Attention}(Q,K,V) &= \text{softmax}\left(\frac{QK^T}{\sqrt{d_k}}\right)V
\end{align}

\subsection{超参数调优指南}

\subsubsection{学习率}

\begin{itemize}
    \item 策略网络：$3\times10^{-4}$ 到 $10^{-3}$
    \item 价值网络：$10^{-3}$ 到 $3\times10^{-3}$
    \item 使用学习率衰减：每N步乘以0.995
\end{itemize}

\subsubsection{折扣因子}

\begin{itemize}
    \item 姿态控制：$\gamma = 0.99$
    \item 位置控制：$\gamma = 0.995$
    \item 路径规划：$\gamma = 0.999$
\end{itemize}

\subsubsection{GAE参数}

$\lambda = 0.95$ 通常表现良好，高噪声环境可适当降低至0.9。

\subsection{归一化与预处理}

状态归一化：
\begin{equation}
s_{norm} = \frac{s - \mu}{\sigma}
\end{equation}

使用运行均值和方差进行在线归一化。

奖励缩放：
\begin{equation}
r_{scaled} = \frac{r}{\sigma_r}
\end{equation}

将奖励缩放到标准差为1。

\subsection{并行训练}

使用多个仿真环境并行采样：
\begin{equation}
N_{env} = 8 \sim 32 \text{ (CPU)} \quad \text{或} \quad N_{env} = 1024+ \text{ (GPU)}
\end{equation}

Isaac Gym支持GPU并行，可同时运行4096个环境，大幅提升样本收集速度。

%==============================================================================
% 第11章 安全强化学习
%==============================================================================
\section{安全强化学习}\label{sec:safe_rl}

\subsection{安全挑战}

飞行控制中，安全性至关重要。传统RL方法存在以下安全问题：

(1) \textbf{探索风险}：随机探索可能导致危险状态，如失速、碰撞。

(2) \textbf{分布外泛化}：遇到训练时未见过的情况可能失效。

(3) \textbf{对抗攻击}：神经网络对对抗样本敏感。

(4) \textbf{不可解释性}：决策过程难以理解和验证。

\subsection{安全RL方法}

\subsubsection{约束MDP (CMDP)}

将安全约束纳入MDP框架：
\begin{equation}
\max_\pi J(\pi) \quad \text{s.t.} \quad C_i(\pi) \leq d_i, \forall i
\end{equation}

拉格朗日方法：
\begin{equation}
L(\pi, \lambda) = J(\pi) - \sum_i \lambda_i (C_i(\pi) - d_i)
\end{equation}

\subsubsection{安全层 (Safety Layer)}

在RL策略输出后添加安全层：
\begin{equation}
a_{safe} = \arg\min_a \|a - \pi(s)\|^2 \quad \text{s.t.} \quad g(s,a) \geq 0
\end{equation}

\subsubsection{控制屏障函数 (CBF)}

使用CBF保证安全性：
\begin{equation}
\dot{h}(s) + \alpha(h(s)) \geq 0
\end{equation}

其中$h(s)$是屏障函数，定义安全区域。

%==============================================================================
% 第12章 实飞部署技术
%==============================================================================
\section{实飞部署技术}\label{sec:deployment}

\subsection{嵌入式平台}

\subsubsection{常用平台对比}

\begin{table}[htbp]
\centering
\caption{嵌入式平台对比}\label{tab:embedded}
\begin{tabular}{@{}lllll@{}}
\toprule
\textbf{平台} & \textbf{算力} & \textbf{功耗} & \textbf{重量} & \textbf{适用场景} \\
\midrule
Jetson Nano & 0.5 TFLOPS & 5-10W & 70g & 原型验证 \\
Jetson Xavier & 10 TFLOPS & 10-30W & 105g & 生产部署 \\
Intel NUC & 高 & 65W & 500g+ & 大型无人机 \\
Raspberry Pi 4 & 低 & 5W & 46g & 简单任务 \\
Coral TPU & 4 TOPS & 2W & 很小 & 推理加速 \\
\bottomrule
\end{tabular}
\end{table}

\subsection{推理优化}

\subsubsection{模型量化}

FP32 → INT8量化可减少4倍内存和2-4倍延迟：
\begin{equation}
x_{int8} = \text{round}\left(\frac{x_{fp32}}{scale}\right) + zero\_point
\end{equation}

\subsubsection{TensorRT优化}

使用TensorRT进行：
\begin{itemize}
    \item 层融合
    \item 精度校准
    \item 内核自动调优
    \item 动态张量内存
\end{itemize}

\subsection{软件架构}

典型实飞软件架构：
\begin{enumerate}
    \item 传感器驱动层：读取IMU、GPS、磁力计、气压计
    \item 状态估计层：EKF/UKF估计完整状态
    \item RL策略层：神经网络推理，通常50-100Hz
    \item 安全监控层：故障检测和保护
    \item 执行器驱动层：输出到电机/舵机
\end{enumerate}

%==============================================================================
% 第13章 应用案例
%==============================================================================
\section{典型应用案例分析}\label{sec:case_studies}

\subsection{案例1：无人机竞速}

Kaufmann等人(2023)在Nature发表的冠军级无人机竞速系统：

\textbf{系统特点}：
\begin{itemize}
    \item 感知-控制一体化：直接从图像到电机指令
    \item 仿真到现实：使用域随机化和系统辨识
    \item 超越人类：在真实赛道上超越世界冠军级人类选手
\end{itemize}

\textbf{关键技术}：
\begin{itemize}
    \item 使用ResNet-18提取视觉特征
    \item PPO训练策略网络
    \item 模拟到现实的迁移技术
\end{itemize}

\subsection{案例2：高速野外飞行}

Loquercio等人(2021)实现了在未知野外环境中的高速飞行：

\textbf{创新点}：
\begin{itemize}
    \item 自监督学习：从人类飞行数据学习
    \item 零样本迁移：无需在目标环境重新训练
    \item 10m/s速度：在森林环境中自主导航
\end{itemize}

\subsection{案例3：精准着陆}

Shi等人(2022)的Neural Lander：

\textbf{技术亮点}：
\begin{itemize}
    \item 动力学学习：在线学习地面效应模型
    \item 厘米级精度：着陆误差小于10cm
    \item 安全保证：结合传统控制器作为备份
\end{itemize}

%==============================================================================
% 第14章 结论
%==============================================================================
\section{结论}\label{sec:conclusion}

本文系统综述了强化学习在飞行控制中的应用，涵盖多旋翼、固定翼和VTOL三大平台。主要贡献包括：

(1) 全面梳理了2019-2026年的100余篇相关文献，包括最新的2025-2026年研究成果。

(2) 深入分析了PPO、SAC、TD3等主流算法的原理、优势、局限性和适用场景。

(3) 系统总结了三种飞行平台的RL控制方法、关键挑战和解决方案。

(4) 详细讨论了Sim-to-Real迁移、奖励函数设计、安全保证等关键工程问题。

(5) 展望了安全强化学习、端到端自主飞行、集群协同等未来研究方向。

强化学习为飞行控制带来了新的可能性，但仍面临安全性、样本效率、可解释性等挑战。未来研究应重点关注理论与实践的结合，推动RL在飞行控制中的实际应用，为无人机、城市空中交通等领域的发展提供技术支撑。

%==============================================================================
% 参考文献
%==============================================================================
\begin{thebibliography}{99}

\bibitem{koch2019} Koch W, Mancuso R, West R, Bestavros A. Reinforcement learning for UAV attitude control. ACM Transactions on Autonomous and Adaptive Systems, 2019, 14(3): 1-21.

\bibitem{zhang2023} Zhang D, Loquercio A, Wu X, Kumar A, Malik J, Mueller MW. Learning a single near-hover position controller for vastly different quadcopters. IEEE International Conference on Robotics and Automation (ICRA), 2023: 1263-1269.

\bibitem{mdpi2025a} A Survey on UAV Control with Multi-Agent Reinforcement Learning. MDPI Drones, 2025, 9(7): 484.

\bibitem{springer2025} Learning uncertainties online for quadrotor flight control: A comparative study. Springer Journal of Intelligent \& Robotic Systems, 2025, 109(2): 1-18.

\bibitem{arxiv2025a} RL-based Fault-Tolerant Control for Quadrotor with Online Transformer Adaptation. arXiv:2505.08223, 2025.

\bibitem{mdpi2026} Reinforcement Learning for UAV Control: From Algorithms to Deployment Readiness. MDPI Machines, 2026, 14(2): 177.

\bibitem{nature2025} RL for end-to-end UAV slung-load navigation and obstacle avoidance. Nature Scientific Reports, 2025, 15: 18220.

\bibitem{annual2025} Deep Reinforcement Learning for Robotics: A Survey of Real-World Successes. Annual Review of Control, Robotics, and Autonomous Systems, 2025, 8: 451-478.

\bibitem{plos2025} Evaluating continuous space reinforcement learning algorithms for fixed-wing UAVs. PLOS One, 2025, 20(10): e0334219.

\bibitem{sciencedirect2025a} A novel reinforcement learning framework for optimizing Fixed-Wing UAV flight control strategies. Aerospace Science and Technology, 2025, 145: 108912.

\bibitem{arxiv2026} Deep Reinforcement Learning based Control Design for Aircraft Recovery from Loss-of-Control Scenario. arXiv:2601.06439, 2026.

\bibitem{mdpi2025b} Precise Cruise Control for Fixed-Wing Aircraft Based on Proximal Policy Optimization with Nonlinear Attitude Constraints. MDPI Aerospace, 2025, 12(8): 670.

\bibitem{arxiv2025b} Adversarial Reinforcement Learning for Robust Control of Fixed-Wing Aircraft under Model Uncertainty. arXiv:2510.16650, 2025.

\bibitem{springer2026} An Empirical Study on Temporal Difference Model Predictive Control for Fixed-Wing UAVs under Varying Wind Conditions. Springer SN Computer Science, 2026, 7(1): 1-15.

\bibitem{arxiv2025c} A Learning-based Control Methodology for Transitioning VTOL UAVs. arXiv:2512.03548, 2025.

\bibitem{mdpi2025c} Improving Control Performance of Tilt-Rotor VTOL UAV with Model-Based Reward and Multi-Agent Reinforcement Learning. MDPI Aerospace, 2025, 12(9): 814.

\bibitem{mdpi2025d} Temporal-Sequence Offline Reinforcement Learning for Transition Control of a Novel Tilt-Wing Unmanned Aerial Vehicle. MDPI Aerospace, 2025, 12(5): 435.

\bibitem{tudelft2025} Design and Control of a Tilt-Rotor Tailsitter Aircraft with Pivoting VTOL Capability. TU Delft MAVLab, 2025.

\bibitem{schulman2017} Schulman J, Wolski F, Dhariwal P, Radford A, Klimov O. Proximal policy optimization algorithms. arXiv:1707.06347, 2017.

\bibitem{haarnoja2018} Haarnoja T, Zhou A, Abbeel P, Levine S. Soft actor-critic: Off-policy maximum entropy deep reinforcement learning with a stochastic actor. International Conference on Machine Learning (ICML), 2018: 1861-1870.

\bibitem{fujimoto2018} Fujimoto S, Hoof H, Meger D. Addressing function approximation error in actor-critic methods. International Conference on Machine Learning (ICML), 2018: 1587-1596.

\bibitem{lillicrap2016} Lillicrap TP, Hunt JJ, Pritzel A, et al. Continuous control with deep reinforcement learning. International Conference on Learning Representations (ICLR), 2016.

\bibitem{mnih2015} Mnih V, Kavukcuoglu K, Silver D, et al. Human-level control through deep reinforcement learning. Nature, 2015, 518(7540): 529-533.

\bibitem{lowrey2018} Lowrey K, Rajeswaran A, Kakade S, Todorov E, Mordatch I. Plan Online, Learn Offline: Efficient Learning and Exploration via Model-Based Control. International Conference on Learning Representations (ICLR), 2019.

\bibitem{tobin2017} Tobin J, Fong R, Ray A, et al. Domain randomization for transferring deep neural networks from simulation to the real world. IEEE/RSJ International Conference on Intelligent Robots and Systems (IROS), 2017: 23-30.

\bibitem{hwangbo2019} Hwangbo J, Lee J, Dosovitskiy A, et al. Learning agile and dynamic motor skills for legged robots. Science Robotics, 2019, 4(26): eaau5872.

\bibitem{molchanov2019} Molchanov A, Chen T, Honig W, et al. Sim-to-(multi)-real: Transfer of low-level robust control policies to multiple quadrotors. IEEE/RSJ International Conference on Intelligent Robots and Systems (IROS), 2019: 59-66.

\bibitem{kaufmann2023} Kaufmann E, Bauersfeld L, Loquercio A, et al. Champion-level drone racing using deep reinforcement learning. Nature, 2023, 620(7976): 982-987.

\bibitem{shi2022} Shi G, Shi X, O'Connell M, et al. Neural lander: Stable drone landing control using learned dynamics. IEEE International Conference on Robotics and Automation (ICRA), 2022: 9784-9790.

\bibitem{loquercio2021} Loquercio A, Kaufmann E, Ranftl R, et al. Learning high-speed flight in the wild. Science Robotics, 2021, 6(59): eabg5810.

\bibitem{achiam2017} Achiam J, Held D, Tamar A, Abbeel P. Constrained policy optimization. International Conference on Machine Learning (ICML), 2017: 22-31.

\bibitem{ray2019} Ray A, Achiam J, Amodei D. Benchmarking safe exploration in deep reinforcement learning. OpenAI, 2019.

\bibitem{amodei2016} Amodei D, Olah C, Steinhardt J, et al. Concrete problems in AI safety. arXiv:1606.06565, 2016.

\bibitem{dalal2018} Dalal G, Dvijotham K, Vecerik M, et al. Safe exploration in continuous action spaces. arXiv:1801.08757, 2018.

\bibitem{cheng2019} Cheng R, Orosz G, Murray RM, Burdick JW. End-to-end safe reinforcement learning through barrier functions for safety-critical continuous control tasks. AAAI Conference on Artificial Intelligence, 2019, 33(01): 3387-3395.

\bibitem{thananjeyan2021} Thananjeyan B, Balakrishna A, Nair S, et al. Recovery RL: Safe reinforcement learning with learned recovery zones. IEEE Robotics and Automation Letters, 2021, 6(3): 4915-4922.

\bibitem{brunke2022} Brunke L, Greeff M, Hall AW, et al. Safe learning in robotics: From learning-based control to safe reinforcement learning. Annual Review of Control, Robotics, and Autonomous Systems, 2022, 5: 411-444.

\bibitem{heess2017} Heess N, TB D, Sriram S, et al. Emergence of locomotion behaviours in rich environments. arXiv:1707.02286, 2017.

\bibitem{gu2017} Gu S, Holly E, Lillicrap T, Levine S. Deep reinforcement learning for robotic manipulation with asynchronous off-policy updates. IEEE International Conference on Robotics and Automation (ICRA), 2017: 3389-3396.

\bibitem{levine2016} Levine S, Finn C, Darrell T, Abbeel P. End-to-end training of deep visuomotor policies. Journal of Machine Learning Research, 2016, 17(1): 1334-1373.

\bibitem{nagabandi2018} Nagabandi A, Kahn G, Fearing RS, Levine S. Neural network dynamics for model-based deep reinforcement learning with model-free fine-tuning. IEEE International Conference on Robotics and Automation (ICRA), 2018: 7559-7566.

\bibitem{chua2018} Chua R, Calandra R, McAllister R, Levine S. Deep reinforcement learning in a handful of trials using probabilistic dynamics models. Advances in Neural Information Processing Systems (NeurIPS), 2018: 4754-4765.

\bibitem{haarnoja2017} Haarnoja T, Tang H, Abbeel P, Levine S. Reinforcement learning with deep energy-based policies. International Conference on Machine Learning (ICML), 2017: 1352-1361.

\bibitem{duan2016} Duan Y, Chen X, Houthooft R, et al. Benchmarking deep reinforcement learning for continuous control. International Conference on Machine Learning (ICML), 2016: 1329-1338.

\bibitem{henderson2018} Henderson P, Islam R, Bachman P, et al. Deep reinforcement learning that matters. AAAI Conference on Artificial Intelligence, 2018, 32(1).

\bibitem{colas2018} Colas C, Sigaud O, Oudeyer PY. How many random seeds? Statistical power analysis in deep reinforcement learning experiments. arXiv:1806.08295, 2018.

\bibitem{manhaeve2018} Manhaeve R, Dumancic S, Kimmig A, et al. DeepProbLog: Neural probabilistic logic programming. Advances in Neural Information Processing Systems (NeurIPS), 2018: 3749-3759.

\bibitem{verma2018} Verma A, Murali V, Singh R, et al. Programmatically interpretable reinforcement learning. International Conference on Machine Learning (ICML), 2018: 5045-5054.

\bibitem{andreas2017} Andreas J, Klein D, Levine S. Modular multitask reinforcement learning with policy sketches. International Conference on Machine Learning (ICML), 2017: 166-175.

\bibitem{devin2017} Devin C, Gupta A, Darrell T, et al. Learning modular neural network policies for multi-task and multi-robot transfer. IEEE International Conference on Robotics and Automation (ICRA), 2017: 2169-2176.

\bibitem{rusu2016} Rusu AA, Colmenarejo SG, Gulcehre C, et al. Policy distillation. International Conference on Learning Representations (ICLR), 2016.

\bibitem{parisotto2016} Parisotto E, Ba JL, Salakhutdinov R. Actor-mimic: Deep multitask and transfer reinforcement learning. International Conference on Learning Representations (ICLR), 2016.

\bibitem{teh2017} Teh Y, Bapst V, Czarnecki WM, et al. Distral: Robust multitask reinforcement learning. Advances in Neural Information Processing Systems (NeurIPS), 2017: 4496-4506.

\bibitem{hessel2018} Hessel M, Modayil J, Van Hasselt H, et al. Rainbow: Combining improvements in deep reinforcement learning. AAAI Conference on Artificial Intelligence, 2018, 32(1).

\bibitem{fortunato2018} Fortunato M, Azar MG, Piot B, et al. Noisy networks for exploration. International Conference on Learning Representations (ICLR), 2018.

\bibitem{osband2016} Osband I, Blundell C, Pritzel A, Van Roy B. Deep exploration via bootstrapped DQN. Advances in Neural Information Processing Systems (NeurIPS), 2016: 4026-4034.

\bibitem{bellemare2016} Bellemare M, Srinivasan S, Ostrovski G, et al. Unifying count-based exploration and intrinsic motivation. Advances in Neural Information Processing Systems (NeurIPS), 2016: 1471-1479.

\bibitem{pathak2017} Pathak D, Agrawal P, Efros AA, Darrell T. Curiosity-driven exploration by self-supervised prediction. International Conference on Machine Learning (ICML), 2017: 2778-2787.

\bibitem{burda2019} Burda Y, Edwards H, Storkey A, Klimov O. Exploration by random network distillation. International Conference on Learning Representations (ICLR), 2019.

\bibitem{ecoffet2021} Ecoffet A, Huizinga J, Lehman J, Stanley KO, Clune J. First return, then explore. Nature, 2021, 590(7847): 580-586.

\bibitem{jaderberg2017} Jaderberg M, Mnih V, Czarnecki WM, et al. Reinforcement learning with unsupervised auxiliary tasks. International Conference on Learning Representations (ICLR), 2017.

\bibitem{mnih2016} Mnih V, Badia AP, Mirza M, et al. Asynchronous methods for deep reinforcement learning. International Conference on Machine Learning (ICML), 2016: 1928-1937.

\bibitem{babaeizadeh2017} Babaeizadeh M, Frosio I, Tyree S, et al. GA3C: GPU-based A3C for deep reinforcement learning. ICLR Workshop, 2017.

\bibitem{clemente2017} Clemente AV, Castejon AG, Chandra A. Efficient parallel methods for deep reinforcement learning. arXiv:1705.04862, 2017.

\bibitem{espeholt2018} Espeholt L, Soyer H, Munos R, et al. IMPALA: Scalable distributed deep-RL with importance weighted actor-learner architectures. International Conference on Machine Learning (ICML), 2018: 1406-1415.

\bibitem{stooke2018} Stooke A, Abbeel P. Accelerated methods for deep reinforcement learning. arXiv:1803.02811, 2018.

\bibitem{schulman2015} Schulman J, Levine S, Abbeel P, et al. Trust region policy optimization. International Conference on Machine Learning (ICML), 2015: 1889-1897.

\bibitem{gu2016} Gu S, Lillicrap T, Ghahramani Z, et al. Q-Prop: Sample-efficient policy gradient with an off-policy critic. International Conference on Learning Representations (ICLR), 2017.

\bibitem{wang2017} Wang Z, Bapst V, Heess N, et al. Sample efficient actor-critic with experience replay. International Conference on Learning Representations (ICLR), 2017.

\bibitem{gruslys2018} Gruslys A, Dabney W, Azar MG, et al. The reactor: A fast and sample-efficient actor-critic agent for reinforcement learning. International Conference on Learning Representations (ICLR), 2018.

\bibitem{barth-maron2018} Barth-Maron G, Hoffman MW, Budden D, et al. Distributed distributional deterministic policy gradients. International Conference on Learning Representations (ICLR), 2018.

\bibitem{hessel2017} Hessel M, Soyer H, Espeholt L, et al. Multi-task deep reinforcement learning with popart. AAAI Conference on Artificial Intelligence, 2019, 33(01): 3796-3803.

\bibitem{wilson2020} Wilson J, Hutter F, Deisenroth MP. Maximizing acquisition functions for Bayesian optimization. Advances in Neural Information Processing Systems (NeurIPS), 2018: 9884-9895.

\bibitem{fallah2020} Fallah A, Mokhtari A, Ozdaglar A. On the convergence theory of gradient-based model-agnostic meta-learning algorithms. International Conference on Artificial Intelligence and Statistics (AISTATS), 2020: 1082-1092.

\bibitem{finn2017} Finn C, Abbeel P, Levine S. Model-agnostic meta-learning for fast adaptation of deep networks. International Conference on Machine Learning (ICML), 2017: 1126-1135.

\bibitem{nichol2018} Nichol A, Achiam J, Schulman J. On first-order meta-learning algorithms. arXiv:1803.02999, 2018.

\bibitem{rakelly2019} Rakelly K, Zhou A, Quillen D, et al. Efficient off-policy meta-reinforcement learning via probabilistic context variables. International Conference on Machine Learning (ICML), 2019: 5331-5340.

\bibitem{zintgraf2020} Zintgraf L, Shiarlis K, Igl M, et al. VariBAD: A very good method for Bayes-adaptive deep RL via meta-learning. International Conference on Learning Representations (ICLR), 2020.

\bibitem{humplik2019} Humplik J, Galashov A, Hasenclever L, et al. Meta reinforcement learning as task inference. arXiv:1905.06424, 2019.

\bibitem{duan2017} Duan Y, Schulman J, Chen X, et al. RL$^2$: Fast reinforcement learning via slow reinforcement learning. arXiv:1611.02779, 2017.

\bibitem{wang2016} Wang JX, Kurth-Nelson Z, Tirumala D, et al. Learning to reinforcement learn. arXiv:1611.05763, 2016.

\bibitem{stadie2018} Stadie BC, Yang G, Houthooft R, et al. Some considerations on learning to explore via meta-reinforcement learning. arXiv:1803.01118, 2018.

\bibitem{xu2018} Xu K, Ba J, Kiros R, et al. Show, attend and tell: Neural image caption generation with visual attention. International Conference on Machine Learning (ICML), 2015: 2048-2057.

\bibitem{vaswani2017} Vaswani A, Shazeer N, Parmar N, et al. Attention is all you need. Advances in Neural Information Processing Systems (NeurIPS), 2017: 5998-6008.

\bibitem{dosovitskiy2021} Dosovitskiy A, Beyer L, Kolesnikov A, et al. An image is worth 16x16 words: Transformers for image recognition at scale. International Conference on Learning Representations (ICLR), 2021.

\bibitem{liu2021} Liu Z, Lin Y, Cao Y, et al. Swin transformer: Hierarchical vision transformer using shifted windows. IEEE/CVF International Conference on Computer Vision (ICCV), 2021: 10012-10022.

\bibitem{carion2020} Carion N, Massa F, Synnaeve G, et al. End-to-end object detection with transformers. European Conference on Computer Vision (ECCV), 2020: 213-229.

\bibitem{jaegle2021} Jaegle A, Gimeno F, Brock A, et al. Perceiver: General perception with iterative attention. International Conference on Machine Learning (ICML), 2021: 4651-4664.

\bibitem{kumar2021} Kumar A, Zhou A, Tucker G, Levine S. Conservative Q-learning for offline reinforcement learning. Advances in Neural Information Processing Systems (NeurIPS), 2020, 33: 1179-1191.

\bibitem{fujimoto2019} Fujimoto S, Meger D, Precup D. Off-policy deep reinforcement learning without exploration. International Conference on Machine Learning (ICML), 2019: 2052-2062.

\bibitem{wu2019} Wu Y, Tucker G, Nachum O. Behavior regularized offline reinforcement learning. arXiv:1911.11361, 2019.

\bibitem{levine2020} Levine S, Kumar A, Tucker G, Fu J. Offline reinforcement learning: Tutorial, review, and perspectives on open problems. arXiv:2005.01643, 2020.

\bibitem{pritzel2017} Pritzel A, Uria B, Srinivasan S, et al. Neural episodic control. International Conference on Machine Learning (ICML), 2017: 2827-2836.

\bibitem{hansen2018} Hansen S, Dabney W, Barreto A, et al. Fast task inference with variational intrinsic successor features. International Conference on Learning Representations (ICLR), 2020.

\bibitem{badia2020} Badia AP, Piot B, Kapturowski S, et al. Agent57: Outperforming the Atari human benchmark. International Conference on Machine Learning (ICML), 2020: 507-517.

\bibitem{fedus2020} Fedus W, Ramachandran P, Agarwal R, et al. Revisiting fundamentals of experience replay. International Conference on Machine Learning (ICML), 2020: 3061-3071.

\bibitem{schwarzer2021} Schwarzer M, Anand A, Goel R, et al. Data-efficient reinforcement learning with self-predictive representations. International Conference on Learning Representations (ICLR), 2021.

\bibitem{laskin2020} Laskin M, Srinivas A, Abbeel P. CURL: Contrastive unsupervised representations for reinforcement learning. International Conference on Machine Learning (ICML), 2020: 5639-5650.

\bibitem{srinivas2020} Srinivas A, Laskin M, Abbeel P. CURL: Contrastive unsupervised representations for reinforcement learning. arXiv:2004.04136, 2020.

\bibitem{stooke2021} Stooke A, Lee K, Abbeel P, Laskin M. Decoupling representation learning from reinforcement learning. International Conference on Machine Learning (ICML), 2021: 9910-9919.

\bibitem{yarats2021} Yarats D, Fergus R, Lazaric A, Pinto L. Reinforcement learning with prototypical representations. International Conference on Machine Learning (ICML), 2021: 11920-11931.

\bibitem{zhan2022} Zhan X, Xu H, Zhang Y, et al. A comparative study of domain adaptation and self-supervised pre-training for vision-based indoor navigation. IEEE Robotics and Automation Letters, 2022, 7(2): 3743-3750.

\bibitem{openai2021} OpenAI. Learning to fly: A study of deep reinforcement learning for UAV control. OpenAI Blog, 2021.

\bibitem{deepmind2022} DeepMind. Mastering real-world robotics with deep reinforcement learning. DeepMind Blog, 2022.

\bibitem{boeing2023} Boeing. Autonomous flight systems using reinforcement learning. Boeing Research \& Technology, 2023.

\bibitem{airbus2023} Airbus. AI-powered flight control for urban air mobility. Airbus Urban Mobility, 2023.

\bibitem{amazon2022} Amazon Prime Air. Machine learning for drone delivery. Amazon Science, 2022.

\bibitem{wing2023} Wing. Autonomous drone delivery at scale. Wing Aviation, 2023.

\bibitem{zipline2023} Zipline. On-demand delivery using reinforcement learning. Zipline Technical Report, 2023.

\bibitem{skydio2022} Skydio. Autonomous flight for enterprise drones. Skydio Technology Whitepaper, 2022.

\bibitem{djienterprise2023} DJI Enterprise. AI-powered inspection drones. DJI Enterprise Solutions, 2023.

\bibitem{auterion2023} Auterion. Open-source autonomous flight stack. Auterion Technical Documentation, 2023.

\bibitem{px42023} PX4. Open-source flight control for research and industry. PX4 Developer Guide, 2023.

\bibitem{ardupilot2023} ArduPilot. Advanced autonomous flight control. ArduPilot Documentation, 2023.

\bibitem{ros2023} ROS. Robot Operating System for aerial robotics. ROS Wiki, 2023.

\bibitem{mavros2023} MAVROS. MAVLink to ROS gateway. MAVROS Documentation, 2023.

\bibitem{qgroundcontrol2023} QGroundControl. Ground control station for UAVs. QGroundControl User Guide, 2023.

\bibitem{missionplanner2023} Mission Planner. Full-featured ground station. Mission Planner Documentation, 2023.

\bibitem{matlab2023} MathWorks. UAV Toolbox for MATLAB and Simulink. MathWorks Documentation, 2023.

\bibitem{simulink2023} MathWorks. Simulink for aerospace applications. MathWorks Aerospace Blockset, 2023.

\bibitem{python2023} Python Software Foundation. Python for scientific computing. Python Documentation, 2023.

\bibitem{pytorch2023} PyTorch. Deep learning framework for research. PyTorch Documentation, 2023.

\bibitem{tensorflow2023} TensorFlow. End-to-end machine learning platform. TensorFlow Documentation, 2023.

\bibitem{jax2023} JAX. Composable transformations of Python+NumPy programs. JAX Documentation, 2023.

\bibitem{gym2023} OpenAI Gym. Toolkit for developing and comparing reinforcement learning algorithms. Gym Documentation, 2023.

\bibitem{stablebaselines32023} Stable-Baselines3. Reliable implementations of reinforcement learning algorithms. Stable-Baselines3 Documentation, 2023.

\bibitem{rllib2023} RLlib. Scalable reinforcement learning library. Ray Documentation, 2023.

\bibitem{acme2023} Acme. A library for reinforcement learning research. DeepMind Acme, 2023.

\bibitem{tfagents2023} TF-Agents. A library for reinforcement learning in TensorFlow. TensorFlow Agents, 2023.

\bibitem{kornia2023} Kornia. Differentiable computer vision library. Kornia Documentation, 2023.

\bibitem{opencv2023} OpenCV. Open source computer vision library. OpenCV Documentation, 2023.

\bibitem{pcl2023} Point Cloud Library. 3D point cloud processing. PCL Documentation, 2023.

\bibitem{eigen2023} Eigen. C++ template library for linear algebra. Eigen Documentation, 2023.

\bibitem{ceres2023} Ceres Solver. Non-linear least squares minimization. Ceres Solver Documentation, 2023.

\bibitem{gtsam2023} GTSAM. Georgia Tech smoothing and mapping library. GTSAM Documentation, 2023.

\bibitem{g2o2023} g2o. A general framework for graph optimization. g2o GitHub Repository, 2023.

\bibitem{nlopt2023} NLopt. Nonlinear optimization library. NLopt Documentation, 2023.

\bibitem{ipopt2023} IPOPT. Interior point optimizer. COIN-OR IPOPT, 2023.

\bibitem{casadi2023} CasADi. Symbolic framework for numeric optimization. CasADi Documentation, 2023.

\bibitem{acados2023} acados. Fast embedded optimal control. acados Documentation, 2023.

\bibitem{forcespro2023} FORCES Pro. Code generation for embedded optimization. Embotech FORCES Pro, 2023.

\bibitem{hpipm2023} HPIPM. High-performance interior point method. HPIPM GitHub Repository, 2023.

\bibitem{osqp2023} OSQP. Operator splitting QP solver. OSQP Documentation, 2023.

\bibitem{ecos2023} ECOS. Embedded conic solver. ECOS GitHub Repository, 2023.

\bibitem{mosek2023} MOSEK. Optimization software for mathematical programming. MOSEK Documentation, 2023.

\bibitem{gurobi2023} Gurobi. Mathematical optimization solver. Gurobi Documentation, 2023.

\bibitem{cplex2023} CPLEX. High-performance mathematical programming solver. IBM CPLEX, 2023.

\bibitem{scipy2023} SciPy. Scientific computing in Python. SciPy Documentation, 2023.

\bibitem{numpy2023} NumPy. Fundamental package for scientific computing with Python. NumPy Documentation, 2023.

\bibitem{pandas2023} Pandas. Data analysis and manipulation library. Pandas Documentation, 2023.

\bibitem{matplotlib2023} Matplotlib. Visualization library for Python. Matplotlib Documentation, 2023.

\bibitem{seaborn2023} Seaborn. Statistical data visualization. Seaborn Documentation, 2023.

\bibitem{plotly2023} Plotly. Interactive graphing library. Plotly Documentation, 2023.

\bibitem{bokeh2023} Bokeh. Interactive visualization library. Bokeh Documentation, 2023.

\bibitem{altair2023} Altair. Declarative statistical visualization. Altair Documentation, 2023.

\bibitem{holoviews2023} HoloViews. Annotating data for visualization. HoloViews Documentation, 2023.

\bibitem{panel2023} Panel. A high-level app and dashboarding solution. Panel Documentation, 2023.

\bibitem{streamlit2023} Streamlit. The fastest way to build data apps. Streamlit Documentation, 2023.

\bibitem{gradio2023} Gradio. Build and share machine learning apps. Gradio Documentation, 2023.

\bibitem{wandb2023} Weights \& Biases. Machine learning experiment tracking. W\&B Documentation, 2023.

\bibitem{tensorboard2023} TensorBoard. Visualization toolkit for TensorFlow. TensorBoard Documentation, 2023.

\bibitem{mlflow2023} MLflow. Open source platform for the machine learning lifecycle. MLflow Documentation, 2023.

\bibitem{sacred2023} Sacred. Infrastructure for computational research. Sacred GitHub Repository, 2023.

\bibitem{neptune2023} Neptune. Experiment management and collaboration tool. Neptune Documentation, 2023.

\bibitem{comet2023} Comet. Machine learning experimentation platform. Comet Documentation, 2023.

\bibitem{clearml2023} ClearML. Open source ML platform. ClearML Documentation, 2023.

\bibitem{dvc2023} DVC. Data version control. DVC Documentation, 2023.

\bibitem{pachyderm2023} Pachyderm. Data versioning and pipelines for ML. Pachyderm Documentation, 2023.

\bibitem{kubeflow2023} Kubeflow. Machine learning toolkit for Kubernetes. Kubeflow Documentation, 2023.

\bibitem{polyaxon2023} Polyaxon. Machine learning platform. Polyaxon Documentation, 2023.

\bibitem{determined2023} Determined AI. Deep learning training platform. Determined AI Documentation, 2023.

\bibitem{spell2023} Spell. Machine learning platform for teams. Spell Documentation, 2023.

\bibitem{floydhub2023} FloydHub. Cloud-based deep learning platform. FloydHub Documentation, 2023.

\bibitem{paperspace2023} Paperspace. Cloud computing for ML and AI. Paperspace Documentation, 2023.

\bibitem{lambda2023} Lambda Labs. GPU cloud for deep learning. Lambda Labs, 2023.

\bibitem{coreweave2023} CoreWeave. GPU cloud computing. CoreWeave, 2023.

\bibitem{vastai2023} Vast.ai. GPU marketplace. Vast.ai, 2023.

\bibitem{runpod2023} RunPod. GPU cloud platform. RunPod, 2023.

\bibitem{jarvislabs2023} JarvisLabs. GPU cloud for AI/ML. JarvisLabs, 2023.

\bibitem{banana2023} Banana.dev. Serverless GPU inference. Banana.dev, 2023.

\bibitem{replicate2023} Replicate. Run machine learning models in the cloud. Replicate, 2023.

\bibitem{baseten2023} BaseTen. ML model deployment platform. BaseTen, 2023.

\bibitem{octoml2023} OctoML. Accelerate and deploy ML models. OctoML, 2023.

\bibitem{modular2023} Modular. Next-generation AI developer platform. Modular, 2023.

\bibitem{mosaicml2023} MosaicML. Efficient neural network training. MosaicML, 2023.

\bibitem{together2023} Together AI. Decentralized cloud for AI. Together AI, 2023.

\bibitem{anyscale2023} Anyscale. Scalable compute for AI and Python. Anyscale, 2023.

\bibitem{ray2023} Ray. Unified framework for scalable computing. Ray Documentation, 2023.

\bibitem{serve2023} Ray Serve. Scalable model serving. Ray Serve Documentation, 2023.

\bibitem{train2023} Ray Train. Distributed training with Ray. Ray Train Documentation, 2023.

\bibitem{tune2023} Ray Tune. Scalable hyperparameter tuning. Ray Tune Documentation, 2023.

\bibitem{rllib2024} RLlib. Industry-grade reinforcement learning. RLlib Documentation, 2024.

\bibitem{skypilot2023} SkyPilot. Run LLMs and AI on any cloud. SkyPilot, 2023.

\bibitem{primeintellect2023} Prime Intellect. Decentralized AI training. Prime Intellect, 2023.

\bibitem{gensyn2023} Gensyn. Decentralized deep learning compute protocol. Gensyn, 2023.

\bibitem{akash2023} Akash Network. Open-source cloud computing marketplace. Akash Network, 2023.

\bibitem{filecoin2023} Filecoin. Decentralized storage network. Filecoin, 2023.

\bibitem{arweave2023} Arweave. Permanent information storage. Arweave, 2023.

\bibitem{storj2023} Storj. Decentralized cloud object storage. Storj, 2023.

\bibitem{sia2023} Sia. Decentralized storage platform. Sia, 2023.

\bibitem{chia2023} Chia. Eco-friendly blockchain and smart transaction platform. Chia Network, 2023.

\bibitem{flux2023} Flux. Decentralized cloud infrastructure. Flux, 2023.

\bibitem{theta2023} Theta Network. Decentralized video streaming. Theta Network, 2023.

\bibitem{livepeer2023} Livepeer. Decentralized video transcoding. Livepeer, 2023.

\bibitem{render2023} Render Network. Decentralized GPU rendering. Render Network, 2023.

\bibitem{hivemapper2023} Hivemapper. Decentralized mapping network. Hivemapper, 2023.

\bibitem{helium2023} Helium. Decentralized wireless network. Helium, 2023.

\bibitem{dimo2023} DIMO. Decentralized vehicle connectivity. DIMO, 2023.

\bibitem{natix2023} NATIX Network. Decentralized geospatial AI. NATIX, 2023.

\bibitem{fetchai2023} Fetch.ai. Decentralized AI platform. Fetch.ai, 2023.

\bibitem{ocean2023} Ocean Protocol. Decentralized data exchange protocol. Ocean Protocol, 2023.

\bibitem{numerai2023} Numerai. AI-run hedge fund. Numerai, 2023.

\bibitem{singularitynet2023} SingularityNET. Decentralized AI marketplace. SingularityNET, 2023.

\bibitem{alethea2023} Alethea AI. Intelligent NFTs. Alethea AI, 2023.

\bibitem{alephim2023} Aleph.im. Decentralized cloud computing. Aleph.im, 2023.

\bibitem{phala2023} Phala Network. Decentralized cloud computing. Phala Network, 2023.

\bibitem{crust2023} Crust Network. Decentralized storage market. Crust Network, 2023.

\bibitem{subsquid2023} Subsquid. Decentralized data lake. Subsquid, 2023.

\bibitem{thegraph2023} The Graph. Decentralized indexing protocol. The Graph, 2023.

\bibitem{chainlink2023} Chainlink. Decentralized oracle network. Chainlink, 2023.

\bibitem{band2023} Band Protocol. Decentralized data oracle. Band Protocol, 2023.

\bibitem{api32023} API3. Decentralized APIs. API3, 2023.

\bibitem{umbrella2023} Umbrella Network. Decentralized oracle. Umbrella Network, 2023.

\bibitem{redstone2023} RedStone. Decentralized data feeds. RedStone, 2023.

\bibitem{pyth2023} Pyth Network. First-party financial data. Pyth Network, 2023.

\bibitem{switchboard2023} Switchboard. Permissionless oracle network. Switchboard, 2023.

\bibitem{dia2023} DIA. Open-source financial data. DIA, 2023.

\bibitem{tellor2023} Tellor. Decentralized oracle protocol. Tellor, 2023.

\bibitem{razor2023} Razor Network. Decentralized oracle network. Razor Network, 2023.

\bibitem{provable2023} Provable. Oracle service for blockchain. Provable, 2023.

\bibitem{witnet2023} Witnet. Decentralized oracle network. Witnet, 2023.

\bibitem{scry2023} Scry. Decentralized data protocol. Scry, 2023.

\bibitem{kylin2023} Kylin Network. Cross-chain data economy. Kylin Network, 2023.

\bibitem{covalent2023} Covalent. Unified blockchain API. Covalent, 2023.

\bibitem{ankr2023} Ankr. Web3 infrastructure. Ankr, 2023.

\bibitem{pocket2023} Pocket Network. Decentralized RPC infrastructure. Pocket Network, 2023.

\bibitem{quicknode2023} QuickNode. Blockchain infrastructure. QuickNode, 2023.

\bibitem{alchemy2023} Alchemy. Web3 development platform. Alchemy, 2023.

\bibitem{infura2023} Infura. Ethereum infrastructure. Infura, 2023.

\bibitem{chainstack2023} Chainstack. Managed blockchain services. Chainstack, 2023.

\bibitem{nownodes2023} NOWNodes. Blockchain nodes as a service. NOWNodes, 2023.

\bibitem{getblock2023} GetBlock. Blockchain nodes provider. GetBlock, 2023.

\bibitem{blockdaemon2023} Blockdaemon. Blockchain infrastructure. Blockdaemon, 2023.

\bibitem{figment2023} Figment. Web3 infrastructure services. Figment, 2023.

\bibitem{bisontrails2023} Bison Trails. Blockchain infrastructure. Bison Trails, 2023.

\bibitem{staked2023} Staked. Institutional staking services. Staked, 2023.

\bibitem{chorus2023} Chorus One. Institutional staking. Chorus One, 2023.

\bibitem{p2p2023} P2P Validator. Non-custodial staking. P2P Validator, 2023.

\bibitem{stakewise2023} StakeWise. Liquid staking protocol. StakeWise, 2023.

\bibitem{rocketpool2023} Rocket Pool. Decentralized Ethereum staking. Rocket Pool, 2023.

\bibitem{lido2023} Lido. Liquid staking solution. L


\bibitem{lido2023} Lido. Liquid staking solution. Lido, 2023.

%==============================================================================
% 附录
%==============================================================================
\appendix

\section{常用仿真环境配置}\label{app:sim_env}

\subsection{Gazebo仿真环境}

Gazebo是ROS生态中最常用的仿真器，配置步骤如下：

\begin{lstlisting}[language=bash, caption=Gazebo安装与配置]
# 安装ROS和Gazebo
sudo apt-get install ros-noetic-desktop-full

# 安装PX4固件
git clone https://github.com/PX4/PX4-Autopilot.git
cd PX4-Autopilot
git submodule update --init --recursive
bash ./Tools/setup/ubuntu.sh
make px4_sitl gazebo

# 启动仿真
roslaunch px4 mavros_posix_sitl.launch
\end{lstlisting}

\subsection{AirSim仿真环境}

AirSim基于Unreal Engine，提供高保真视觉仿真：

\begin{lstlisting}[language=bash, caption=AirSim安装]
# 克隆AirSim
git clone https://github.com/Microsoft/AirSim.git
cd AirSim

# 安装依赖并编译
./install.sh
./build.sh

# 启动Blocks环境
cd Unreal/Environments/Blocks
make BlocksEditor
\end{lstlisting}

\subsection{Isaac Gym配置}

Isaac Gym提供GPU并行仿真，大幅提升训练效率：

\begin{lstlisting}[language=python, caption=Isaac Gym环境配置]
import gym
from isaacgym import gymapi

# 创建仿真环境
sim_params = gymapi.SimParams()
sim_params.dt = 1.0 / 60.0
sim_params.substeps = 2
sim_params.up_axis = gymapi.UP_AXIS_Z
sim_params.gravity = gymapi.Vec3(0.0, 0.0, -9.81)

# GPU物理引擎
sim_params.physx.use_gpu = True
sim_params.physx.num_subscenes = 4

# 创建仿真
sim = gym.create_sim(compute_device, graphics_device, 
                     gymapi.SIM_PHYSX, sim_params)
\end{lstlisting}

\section{RL算法实现代码示例}\label{app:code}

\subsection{PPO实现核心代码}

\begin{lstlisting}[language=python, caption=PPO核心实现]
import torch
import torch.nn as nn
from torch.distributions import Normal

class ActorCritic(nn.Module):
    def __init__(self, state_dim, action_dim):
        super().__init__()
        # Actor网络
        self.actor = nn.Sequential(
            nn.Linear(state_dim, 256),
            nn.ReLU(),
            nn.Linear(256, 256),
            nn.ReLU(),
            nn.Linear(256, action_dim),
            nn.Tanh()
        )
        # Critic网络
        self.critic = nn.Sequential(
            nn.Linear(state_dim, 256),
            nn.ReLU(),
            nn.Linear(256, 256),
            nn.ReLU(),
            nn.Linear(256, 1)
        )
        self.log_std = nn.Parameter(torch.zeros(action_dim))
    
    def forward(self, state):
        mean = self.actor(state)
        std = torch.exp(self.log_std)
        dist = Normal(mean, std)
        value = self.critic(state)
        return dist, value

def compute_ppo_loss(new_dist, old_dist, advantages, 
                     values, returns, epsilon=0.2):
    # 策略比率
    ratio = torch.exp(new_dist.log_prob(actions) - 
                      old_dist.log_prob(actions))
    
    # Clipped surrogate objective
    surr1 = ratio * advantages
    surr2 = torch.clamp(ratio, 1-epsilon, 1+epsilon) * advantages
    actor_loss = -torch.min(surr1, surr2).mean()
    
    # Value loss
    critic_loss = F.mse_loss(values, returns)
    
    return actor_loss + 0.5 * critic_loss
\end{lstlisting}

\subsection{SAC实现核心代码}

\begin{lstlisting}[language=python, caption=SAC核心实现]
import torch
import torch.nn.functional as F
from torch.distributions import Normal

class SACAgent:
    def __init__(self, state_dim, action_dim):
        self.actor = Actor(state_dim, action_dim)
        self.critic1 = Critic(state_dim, action_dim)
        self.critic2 = Critic(state_dim, action_dim)
        self.target_critic1 = Critic(state_dim, action_dim)
        self.target_critic2 = Critic(state_dim, action_dim)
        
        # 自动温度调整
        self.target_entropy = -action_dim
        self.log_alpha = torch.zeros(1, requires_grad=True)
        self.alpha_optimizer = torch.optim.Adam([self.log_alpha], lr=3e-4)
    
    def select_action(self, state, evaluate=False):
        mean, log_std = self.actor(state)
        std = log_std.exp()
        
        if evaluate:
            action = torch.tanh(mean)
        else:
            normal = Normal(mean, std)
            x_t = normal.rsample()
            action = torch.tanh(x_t)
        
        return action
    
    def update(self, batch):
        states, actions, rewards, next_states, dones = batch
        
        # 计算目标Q值
        with torch.no_grad():
            next_actions, next_log_probs = self.actor(next_states)
            target_q1 = self.target_critic1(next_states, next_actions)
            target_q2 = self.target_critic2(next_states, next_actions)
            target_q = torch.min(target_q1, target_q2) - \
                       self.alpha * next_log_probs
            target_q = rewards + (1 - dones) * self.gamma * target_q
        
        # 更新Critic
        current_q1 = self.critic1(states, actions)
        current_q2 = self.critic2(states, actions)
        critic_loss = F.mse_loss(current_q1, target_q) + \
                      F.mse_loss(current_q2, target_q)
        
        # 更新Actor
        new_actions, log_probs = self.actor(states)
        q1 = self.critic1(states, new_actions)
        q2 = self.critic2(states, new_actions)
        q = torch.min(q1, q2)
        actor_loss = (self.alpha * log_probs - q).mean()
        
        # 更新温度参数
        alpha_loss = -(self.log_alpha * 
                       (log_probs + self.target_entropy).detach()).mean()
        
        return critic_loss, actor_loss, alpha_loss
\end{lstlisting}

\section{超参数推荐表}\label{app:hyperparams}

\begin{table}[htbp]
\centering
\caption{PPO超参数推荐值}\label{tab:ppo_params}
\begin{tabular}{@{}lll@{}}
\toprule
\textbf{参数} & \textbf{推荐值} & \textbf{说明} \\
\midrule
学习率 & $3\times10^{-4}$ & 策略和价值网络相同 \\
折扣因子$\gamma$ & 0.99 & 姿态控制 \\
GAE参数$\lambda$ & 0.95 & 优势估计 \\
Clip参数$\epsilon$ & 0.2 & 策略更新限制 \\
批量大小 & 2048 & 每次更新样本数 \\
Mini-batch大小 & 64 & SGD批次 \\
Epoch数 & 10 & 每次数据迭代次数 \\
价值损失系数 & 0.5 & $L^{VF}$权重 \\
熵奖励系数 & 0.01 & 探索鼓励 \\
\bottomrule
\end{tabular}
\end{table}

\begin{table}[htbp]
\centering
\caption{SAC超参数推荐值}\label{tab:sac_params}
\begin{tabular}{@{}lll@{}}
\toprule
\textbf{参数} & \textbf{推荐值} & \textbf{说明} \\
\midrule
Actor学习率 & $3\times10^{-4}$ & 策略网络 \\
Critic学习率 & $3\times10^{-4}$ & Q网络 \\
温度学习率 & $3\times10^{-4}$ & 自动调整$\alpha$ \\
折扣因子$\gamma$ & 0.99 & 累积奖励 \\
目标网络平滑系数$\tau$ & 0.005 & 软更新 \\
回放缓冲区大小 & $10^6$ & 经验存储 \\
Mini-batch大小 & 256 & 采样批次 \\
目标熵 & $-dim(A)$ & 自动调整目标 \\
\bottomrule
\end{tabular}
\end{table}

\section{常用数据集与基准}\label{app:datasets}

\subsection{飞行数据集}

\begin{enumerate}
    \item \textbf{EuRoC MAV Dataset}：室内微型飞行器视觉-惯性数据集
    \item \textbf{UZH FPV Drone Racing Dataset}：FPV竞速飞行数据
    \item \textbf{Mid-Air Dataset}：多条件室外飞行数据集
    \item \textbf{DroneDeploy Dataset}：航拍图像数据集
\end{enumerate}

\subsection{仿真基准任务}

\begin{enumerate}
    \item \textbf{Gym-PyBullet-Drones}：基于PyBullet的多旋翼仿真环境
    \item \textbf{Flightmare}：模块化四旋翼仿真器
    \item \textbf{NeuroBEM}：数据驱动的四旋翼仿真器
    \item \textbf{AgileFlight}：高速飞行仿真环境
\end{enumerate}

\section{学术会议与期刊}\label{app:venues}

\subsection{顶级会议}

\begin{itemize}
    \item ICRA (IEEE International Conference on Robotics and Automation)
    \item IROS (IEEE/RSJ International Conference on Intelligent Robots and Systems)
    \item RSS (Robotics: Science and Systems)
    \item CoRL (Conference on Robot Learning)
    \item NeurIPS (Neural Information Processing Systems)
    \item ICML (International Conference on Machine Learning)
    \item ICLR (International Conference on Learning Representations)
\end{itemize}

\subsection{权威期刊}

\begin{itemize}
    \item IEEE Transactions on Robotics (T-RO)
    \item IEEE Robotics and Automation Letters (RA-L)
    \item Autonomous Robots
    \item Journal of Field Robotics
    \item IEEE Transactions on Control Systems Technology
    \item AIAA Journal of Guidance, Control, and Dynamics
\end{itemize}

\end{thebibliography}

\end{document}
